\section{Orígenes y evolución histórica de la ingeniería de software}

La ingeniería de software surgió como respuesta a la necesidad de transformar la programación, que durante sus primeras etapas fue una actividad predominantemente artesanal, en una disciplina capaz de producir sistemas complejos con mayor control, calidad y previsibilidad. En los inicios de la computación, especialmente en las décadas de 1940 y 1950, el desarrollo de programas se realizaba con poca estandarización, escasa documentación y casi sin métodos formales de diseño, lo que hacía muy difícil el mantenimiento y la evolución de los sistemas \cite{sommerville2016}.

Con el crecimiento del tamaño y la complejidad de los sistemas informáticos, comenzaron a aparecer numerosos problemas de retrasos, sobrecostos, fallos de fiabilidad y dificultad para adaptar el software a nuevas necesidades. Esta situación dio lugar a lo que posteriormente se denominó la crisis del software, un fenómeno ampliamente reconocido desde finales de los años 60 y consolidado en la Conferencia de la NATO sobre Ingeniería de Software de 1968, donde se propuso tratar el desarrollo de software como una verdadera disciplina de ingeniería \cite{naur1969}. A partir de ese momento, el término ``ingeniería de software'' comenzó a difundirse para designar un enfoque sistemático y disciplinado de construcción de programas \cite{boehm1981}.

\subsection{Décadas de 1960 y 1970: primeros modelos de ciclo de vida}

Durante las décadas de 1960 y 1970, el aumento de la complejidad de los sistemas impulsó la búsqueda de modelos de desarrollo más estructurados. En este contexto apareció el modelo en cascada, una propuesta secuencial que divide el proceso en fases bien definidas, como análisis de requisitos, diseño, implementación, pruebas y mantenimiento \cite{royce1970}. Aunque más tarde fue criticado por su rigidez, el modelo en cascada fue importante porque introdujo una visión ordenada del ciclo de vida del software y reforzó la idea de que cada etapa debía documentarse y validarse antes de avanzar a la siguiente \cite{pressman2020}.

En estos años también se fortaleció el énfasis en la documentación, las pruebas y la modularidad. La documentación se consideró esencial para facilitar la comunicación entre desarrolladores, usuarios y gestores, así como para asegurar la mantenibilidad de los sistemas \cite{sommerville2016}. Las pruebas comenzaron a consolidarse como una actividad indispensable para detectar errores antes de la entrega, y la modularidad se presentó como una estrategia para dividir sistemas grandes en componentes más manejables y reutilizables \cite{parnas1972}.

\subsection{Décadas de 1980 y 1990: profesionalización y calidad}

A partir de los años 80, la ingeniería de software entró en una etapa de mayor profesionalización. El crecimiento de la industria del software exigió prácticas más maduras de planificación, gestión, control de calidad y mejora continua. Uno de los marcos más influyentes fue el Capability Maturity Model (CMM), desarrollado en el Software Engineering Institute, que propuso niveles de madurez para evaluar y perfeccionar los procesos de desarrollo \cite{paulk1993}. Este modelo influyó profundamente en la manera en que las organizaciones comenzaron a ver la calidad no solo como una propiedad del producto final, sino también del proceso que lo produce \cite{humphrey1989}.

En paralelo, se consolidaron estándares internacionales como los de la familia ISO/IEC, orientados a establecer criterios comunes para la calidad del software y de los procesos asociados \cite{iso12207}. Estos estándares ayudaron a formalizar prácticas de desarrollo, prueba, mantenimiento y gestión de la configuración, favoreciendo una visión más industrial del software \cite{sommerville2016}. Durante este período también se expandió la Programación Orientada a Objetos (POO), cuyos principios de encapsulación, herencia y polimorfismo facilitaron la construcción de sistemas más flexibles y reutilizables \cite{booch1994}. Asimismo, aparecieron las herramientas CASE (Computer-Aided Software Engineering), diseñadas para apoyar de forma automatizada tareas de análisis, diseño, documentación y mantenimiento \cite{pressman2020}.

\subsection{Décadas de 2000 y 2010: agilidad y automatización}

En los años 2000 y 2010, la ingeniería de software experimentó una transformación importante debido a la necesidad de adaptarse a entornos de cambio rápido. Frente a los modelos pesados y rígidos, cobraron protagonismo las metodologías ágiles, que proponen un desarrollo iterativo, incremental y centrado en la colaboración continua con el cliente \cite{beck2001}. Entre los marcos más influyentes destacan Scrum y Extreme Programming (XP), ambos orientados a obtener retroalimentación frecuente, responder con rapidez a cambios en los requisitos y entregar valor en ciclos cortos \cite{schwaber2020, beck2004}.

En esta misma etapa surgieron prácticas que transformaron la producción y operación del software, especialmente DevOps, que promueve la integración entre desarrollo y operaciones para acelerar la entrega y mejorar la calidad mediante automatización y monitoreo continuo \cite{kim2016}. La integración continua permitió validar de forma frecuente los cambios en el código mediante pruebas automáticas, reduciendo errores e incrementando la confiabilidad \cite{fowler2006}. Posteriormente, la infraestructura como código amplió esta lógica al terreno de la administración de sistemas, permitiendo definir y reproducir entornos de despliegue mediante archivos versionados y automatizados \cite{humble2010}.

\subsection{Síntesis histórica}

En síntesis, la evolución de la ingeniería de software refleja el paso desde una práctica artesanal hacia una disciplina madura, apoyada en métodos, estándares, herramientas y procesos cada vez más sofisticados. Cada etapa histórica respondió a desafíos concretos de su tiempo: primero, la necesidad de organizar la programación; después, la de garantizar calidad y mantenimiento; más tarde, la de mejorar la adaptabilidad; y finalmente, la de automatizar la entrega y operación de sistemas complejos \cite{sommerville2016, pressman2020}. Por ello, la ingeniería de software constituye hoy una de las bases fundamentales de la computación moderna \cite{boehm1981}.

\begin{figure}[H]
    \centering
    % Descomentar y usar tu imagen: \includegraphics[width=0.6\textwidth]{ruta/a/imagen.png}
    \caption{Línea de tiempo de la evolución de la Ingeniería de Software}
    \label{fig:evolucion}
\end{figure}

\noindent\textbf{Versión interactiva:} \\
La línea de tiempo interactiva se encuentra en el archivo HTML incluido en la carpeta del proyecto (.rar). \\
\href{run:timeline.html}{Abrir línea de tiempo interactiva}